\chapter{A complete Gaussian field laboratory}\label{ch:lab}
This laboratory is a sequence of exact worked calculations. It is not a
trajectory produced by the research engine and not a proposed configuration
for a scientific campaign. Its purpose is to let the reader reconstruct
every important distinction with pencil-and-paper arithmetic.

\section{Write the declaration before calculating}
For the first cases, use two cells with
\begin{equation}
 z=(18,2),\quad x^*=(10,10),\quad\sigma=(2,2),
 \quad H=\operatorname{diag}(1/4,1/4).
\end{equation}
Both state and reference sum to twenty. The action is a lossless transfer
from the first cell to the second, with finite quantity $q$. The increment
is $(-q,q)$. The illustrative physical domain permits $0\leq q\leq18$;
there is no separate destination cap in this example. Process burden is zero.

The initial potential is
\begin{equation}
 V(z)=\tfrac12[(8/2)^2+(-8/2)^2]=16.
\end{equation}
The marginals are $(2,-2)$, so the field force is four. Substitution into
the exact finite expression gives $E(q)=4q-q^2/4$. Every later number
should be checked both from that polynomial and from the endpoint potential.

\section{A menu with favorable and adverse quantities}
\begin{center}
\begin{tabular}{rrrr}
\toprule $q$ & Endpoint & $V$ after & $E(q)$\\\midrule
0 & $(18,2)$ & 16 & 0\\
1 & $(17,3)$ & $49/4$ & $15/4$\\
2 & $(16,4)$ & 9 & 7\\
4 & $(14,6)$ & 4 & 12\\
8 & $(10,10)$ & 0 & 16\\
12 & $(6,14)$ & 4 & 12\\
16 & $(2,18)$ & 16 & 0\\
17 & $(1,19)$ & $81/4$ & $-17/4$\\
18 & $(0,20)$ & 25 & $-9$\\\bottomrule
\end{tabular}
\end{center}
Each row conserves twenty and has nonnegative stocks. Therefore physical
feasibility alone does not choose the helpful rows. Conversely, the two
zero-valued quantities mean very different things: $q=0$ does nothing,
while $q=16$ exchanges the deviations and returns to the same potential.
Equal potential is not equal state.

The table is not a controller instruction. Selecting eight because it has
the largest value would define an optimizer. Choosing uniformly from a
declared affordable subset would define another process. Listing analytical
options does not authorize either one as the scientific study.

\section{Check the integral independently}
Along the transfer, $f(r)=4-r/2$. For $q=4$,
\begin{equation}
 \int_0^4(4-r/2)\,dr=[4r-r^2/4]_0^4=12.
\end{equation}
The initial-tangent estimate is sixteen, overvaluing the action by four.
The average force along the segment is $(4+2)/2=3$, and multiplying by
quantity four again gives twelve. The endpoint-force averaging is exact
because this particular gradient is affine, not because trapezoidal
quadrature is exact for all field models.

At $q=12$, the final force is minus two. The average is one and the
integral is twelve. Some portions of the path improve the potential and
later portions worsen it, yet the total remains positive. A positive
integral is not a statement that every infinitesimal part was beneficial.

\section{Separate force from an Onsager response}
Declare, only for this illustration, $M_e=1/2$ and $\theta_e=0$.
At the original state the requested rate is two. Duration $1/2$ gives
quantity one and exact field value $15/4$. Duration two gives quantity
four and value twelve. Duration four gives quantity eight and value sixteen.
These are three frozen-step calculations, not repeated steps in one run.

If instead the flux were recalculated continuously as the state moved,
the force would fall and so would the flux. It would be a mistake to
interpret the frozen duration-four jump as the solution of that differential
equation. ``Same starting force'' does not mean ``same time evolution.''

\section{Local evaluation with a nonzero remote term}
Append a third cell at stock twelve, reference ten and scale two, but do
not change it. The initial global potential becomes $33/2$. After the
four-unit transfer it is $9/2$. The remote half-unit contribution is
present in both and cancels. Local valuation still gives twelve.

Now imagine the transfer also consumes stock from that third cell. The
former two-cell formula is no longer a complete local account. The example
does not need a numerical error to demonstrate the problem: its support
assumption has changed. A correct implementation should reject the incomplete
declaration rather than silently use the old result.

\section{Sequential receipts and reversed order}
Divide a four-unit transfer into quantities one and three. If one occurs
first, its value is $15/4$ and the subsequent three-unit value is
$49/4-4=33/4$. If three occurs first, its value is $39/4$ and the last
one-unit value is $25/4-4=9/4$. Both sums are twelve.

The intermediate state explains the different child values. There is no
contradiction to reconcile by averaging the orders unless a separate
allocation rule asks for such an average. The physical sequence already
defines which live state each action used.

If an extra nonnegative activation burden is incurred for each separate
action, splitting can change the net total. The field sum still telescopes;
the process-cost sum changes. An implementation must not quote the one-action
cost and then execute two independently activated operations under it.

\section{A simultaneous group on the same source}
Let both two-unit children be accepted on the same baseline. The net
increment is $(-4,4)$ and the group field value is twelve. Same-baseline
singletons each show seven, yielding a misleading fourteen if added as
though independent. The missing cross correction is minus two.

Under the declared straight common path, the attribution to each child is
six. This is easy to check from the average group gradient, or by giving
each child half of the identical net direction. Closure is $6+6=12$.
It does not prove that each actor owns six units of spendable capacity.

For unequal children of one and three units, common-path values become
three and nine. These differ from both ordered sequential splits above.
The group field result is the same because the endpoint is the same;
the child interpretations refer to different declared constructions.

\section{Cancellation and a signed child account}
Use the same baseline but take opposing children $(-2,2)$ and $(2,-2)$.
The net endpoint is unchanged. Standalone values are $7$ and $-9$,
while common-path values are $8$ and $-8$. Both the group value and the
sum of path values are zero. The singleton sum is minus two.

If each child belongs to a different zero-balance owner, the proposed
owner-wise affordability rule refuses the owner charged eight. If both
belong to one owner, same-event netting gives zero. This is an explicit
demonstration that ownership policy can matter even when the physical net
change is identical. It is not a recommendation to choose the ownership
map that produces a desired outcome.

\section{A different scale is a different model}
Keep the original state and reference but change scales to $(2,4)$.
The initial potential is ten, marginals are $(2,-1/2)$, and field is
$5/2$. For quantity four the curvature coefficient is
$\tfrac12(1/4+1/16)=5/32$, so
\begin{equation}
 E(4)=4(5/2)-(5/32)16=15/2.
\end{equation}
The endpoint potential is $5/2$, agreeing with $10-5/2=15/2$.
Nothing physical changed between this declaration and the earlier example.
The change in value comes from the declared scale at the second cell.
Any empirical use must therefore justify and version that declaration.

Changing to smaller physical units would instead rescale stock, reference,
quantity and scale together and leave EBU unchanged. A field account that
changes when only the unit label changes has a dimensional bug.

\section{A route whose first step is unfavorable}
Declare a new three-cell example with state $(10,10,2)$, reference
$(6,6,10)$ and equal scales two. Both totals are twenty-two. The initial
potential is twelve. A four-unit transfer from the first cell to the
second gives $(6,14,2)$ with potential sixteen, so its value is minus four.
A subsequent eight-unit transfer from the second cell to the third reaches
the reference $(6,6,10)$ and has value sixteen. The route total is twelve.

This is a physically feasible quantity-fixed sequence under simple stock
nonnegativity. It illustrates the difference between favorable endpoint
value and intermediate affordability. A zero-balance owner cannot pay the
first negative action under the proposed no-borrowing rule. The later
positive action does not retroactively authorize it. A different grouping,
timing or ownership policy would be a different declared experiment.

\section{Forcing and the complete account}
Begin instead at $(10,10)$ with $V=0$, zero balances and zero external
audit. A conservative shock gives $(14,6)$, with $D_{\mathrm{ext}}=4$.
One two-unit restoring transfer ends at $(12,8)$, with potential one and
field value three. If the candidate account policy credits that action
to its declared owner, the invariant reads
\begin{equation}
 1+3-4=0.
\end{equation}
Now consider a possible reverse two-unit action from $(12,8)$ back to
$(14,6)$. Its value is minus three. An owner with balance three can afford
it, returning to potential four and balance zero, still with audit four.
Both accounts close. Recovery has been undone. This is a counterexample
to the claim that the accounting invariant alone forces monotone recovery.

No random process has been run to obtain this pair of possibilities. It is
an algebraic construction showing what the proposed rules allow. A later
study must determine how often such events arise under its actual menus and
sampling, without changing those rules after seeing results.

\section{A reproducible hand calculation}
For every candidate, a reader should be able to list the baseline, references,
scales, child increments, complete support, physical constraints, before/after
potential, group value and optional path values. If capacity is discussed,
the owner map and prior balances must also be listed. If forcing occurs,
the before-forcing state and signed audit must be listed separately.

An unexplained number in any one of those positions should interrupt the
claim. The appropriate response is not to fill it with a plausible value.
It is to identify which declaration, measurement or authority is missing.
This is how a careful laboratory connects to the fail-closed software
architecture described in Part III.

\section*{Final synthesis}
The Gaussian model gives exact, local finite accounting and a transparent
curvature correction. It supports conditional path and telescoping proofs,
and it makes group interaction calculable. The Onsager-type controller gives
an additional conditional dissipation law when it is explicitly chosen.
Random-affordable action is a different prospective process. Its useful,
adverse and inconclusive outcomes remain to be investigated, not written
into the definition of success.
